\documentclass[a4paper, serif, 10pt]{moderncv}

%% moderncv
% style and color
\moderncvstyle{banking}
\moderncvcolor{black}

%% page layout/margins
\usepackage[top=2.5cm, bottom=2.5cm, left=1.5cm, right=1.5cm]{geometry}

\nopagenumbers{}

%% Babel config for German language
% ref:
% - https://latex3.github.io/babel/guides/locale-german.html
\usepackage[ngerman]{babel}

%% fontspec
\usepackage{fontspec}

\IfFontExistsTF{Linux Libertine}{
  \setmainfont[Ligatures={TeX, Common}, Numbers=OldStyle]{Linux Libertine}
}{}
\IfFontExistsTF{Linux Libertine O}{
  \setmainfont[Ligatures={TeX, Common}, Numbers=OldStyle]{Linux Libertine O}
}{}

%% CTeX
% CJK support, used to show CJK characters in the resume
%
% - fontset=none: disable builtin fontset but instead set the CJK font manually
% - heading=false: disable ctex heading
% - punct=kaiming: use kaiming punctuations styles for CJK
% - scheme=plain: use plain scheme, do not override \normalsize font size
% - space=auto: space settings for CJK characters
%
% ref:
% - http://ctan.mirrorcatalogs.com/language/chinese/ctex/ctex.pdf
\usepackage[UTF8, heading=false, punct=kaiming, scheme=plain, space=auto]{ctex}

\IfFontExistsTF{Noto Serif CJK SC}{
  \setCJKmainfont{Noto Serif CJK SC}
}{}
\IfFontExistsTF{Noto Sans CJK SC}{
  \setCJKsansfont{Noto Sans CJK SC}
}{}

%% line spacing
\usepackage{setspace}
\setstretch{1.125}

%% URL styling - use normal text instead of monospace
\urlstyle{same}

% Auto-underline all links
% Defer until \begin{document} so hyperref is already loaded
\AtBeginDocument{%
  \let\oldhref\href
  \renewcommand{\href}[2]{\oldhref{#1}{\underline{#2}}}%
}

%% Patch \httplink and \httpslink to handle full URLs
% The original moderncv macros always prepend http:// or https:// to URLs.
% This causes issues when the URL already has a protocol (e.g., https://example.com)
% resulting in malformed URLs like https://https://example.com.
% This patch checks if the URL already contains "://" and uses it directly if so.
% Note: We use \str_set:Nx (x-expansion) to expand macros like \@homepage first.
\ExplSyntaxOn
\str_new:N \g_yamlresume_url_str

\RenewDocumentCommand{\httpslink}{O{}m}{
  \str_set:Nx \g_yamlresume_url_str {#2}
  \str_if_in:NnTF \g_yamlresume_url_str {://}
    { \url{#2} }
    { \url{https://#2} }
}

\RenewDocumentCommand{\httplink}{O{}m}{
  \str_set:Nx \g_yamlresume_url_str {#2}
  \str_if_in:NnTF \g_yamlresume_url_str {://}
    { \url{#2} }
    { \url{http://#2} }
}
\ExplSyntaxOff

\name{Lena Hoffmann}{}
\title{Data Analystin für Business Intelligence und Datenvisualisierung}
\phone[mobile]{+49 30 12345678}
\email{lena.hoffmann@posteo.de}
\homepage{https://lena-hoffmann.dev}

\address{Berlin, Berlin, Deutschland, 10115}{}{}

\extrainfo{{\small \faLinkedin}\ \href{https://www.linkedin.com/in/lena-hoffmann}{@lena-hoffmann} {} {} {} • {} {} {}
{\small \faGithub}\ \href{https://github.com/lena-hoffmann}{@lena-hoffmann}}

\begin{document}

\maketitle

\section{Basisinformationen}

\cvline{}{Data Analystin mit mehr als vier Jahren Erfahrung in der Analyse von Umsatz-, Kunden- und Produktdaten. Ich entwickle verständliche Dashboards, automatisiere Auswertungen und übersetze Ergebnisse in konkrete Handlungsempfehlungen für Fachbereiche und Management.
%
\begin{itemize}
\item Schwerpunkt: Business Intelligence, Datenvisualisierung und statistische Analyse
\item Werkzeuge: SQL, Python, Power BI, Tableau, dbt
\item Methoden: A/B-Tests, Kohortenanalysen, Regressionsmodelle
\item Stärken: strukturierte Arbeitsweise, Kommunikation auf Augenhöhe und hohe Zahlenaffinität
\end{itemize}}
\section{Ausbildung}

\cventry{Okt. 2016–Sept. 2020}
        {Bachelor, Statistik und Wirtschaftsinformatik, Punkte: 1.7}
        {Humboldt-Universität zu Berlin}
        {\url{https://www.hu-berlin.de}}
        {}
        {\begin{itemize}
\item Schwerpunkt auf statistischer Datenanalyse und praktischer Anwendung von BI-Methoden
\item Bachelorarbeit zum Thema „Prognosemodelle für Umsatzdaten im E-Commerce“ mit Note 1,3
\item Wahlpflichtmodule in Data Mining und multivariater Statistik
\end{itemize}
\textbf{Kurse}: Einführung in die Statistik, Datenbankmanagementsysteme, Betriebswirtschaftslehre, Data Mining, Multivariate Verfahren, Informationsvisualisierung}
\section{Berufserfahrung}

\cventry{März 2022–Heute}
        {Data Analystin}
        {Nordlicht Analytics GmbH}
        {\url{https://www.nordlicht-analytics.de}}
        {}
        {\begin{itemize}
\item Entwicklung und Pflege interaktiver Dashboards und standardisierter Reportings in Power BI für Geschäftsführung und Fachabteilungen
\item Aufbau und Optimierung eines relationalen Datenmodells in PostgreSQL sowie Automatisierung von ETL-Prozessen mit Python und dbt
\item Durchführung von A/B-Tests und Kundensegmentanalysen zur Bewertung von Marketingkampagnen; Identifikation konkreter Optimierungspotenziale
\item Reduzierung manueller Reporting-Tätigkeiten um rund 40 \% durch automatisierte Datenpipelines
\item Enge Zusammenarbeit mit Produktmanagement, Vertrieb und Marketing zur Ableitung datenbasierter Handlungsempfehlungen
\end{itemize}
\textbf{Schlüsselwörter}: Power BI, SQL, Python, dbt, ETL, A/B-Tests}

\cventry{Sept. 2020–Feb. 2022}
        {Junior Data Analystin}
        {Datenwerk Solutions AG}
        {\url{https://www.datenwerk-solutions.de}}
        {}
        {\begin{itemize}
\item Erstellung wöchentlicher Absatz-, Bestands- und Umsatzberichte für das Controlling
\item Bereinigung und Konsolidierung von Stamm- und Bewegungsdaten in Excel und SQL
\item Unterstützung bei der Einführung von Tableau und der Migration bestehender Reports
\item Automatisierung von Routineauswertungen mit Python-Skripten
\item Beantwortung von Ad-hoc-Analysen für Fachbereiche mit präzisen und nachvollziehbaren Ergebnissen
\end{itemize}
\textbf{Schlüsselwörter}: Excel, SQL, Tableau, Python, Berichtswesen}

\cventry{Apr. 2019–Aug. 2020}
        {Studentische Hilfskraft für Datenanalyse}
        {Humboldt-Universität zu Berlin}
        {\url{https://www.hu-berlin.de}}
        {}
        {\begin{itemize}
\item Unterstützung von Forschungsprojekten bei Datenaufbereitung und statistischer Auswertung mit R
\item Durchführung von Regressions- und Varianzanalysen inklusive Dokumentation der Ergebnisse
\item Betreuung des Statistiklabors und Durchführung von Tutorien für Studierende
\end{itemize}
\textbf{Schlüsselwörter}: R, Statistik, SPSS}
\section{Sprachen}

\cvline{Deutsch}{Muttersprache oder zweisprachig \hfill \textbf{Schlüsselwörter}: Muttersprache}
\cvline{Englisch}{Verhandlungssichere Kenntnisse \hfill \textbf{Schlüsselwörter}: TOEFL 105, Fachliche Kommunikation}
\section{Fähigkeiten}

\cvline{Datenanalyse und Statistik}{Experte \hfill \textbf{Schlüsselwörter}: Python, R, SQL, A/B-Tests, Regressionsanalyse, Statistik}
\cvline{Business Intelligence}{Erfahren \hfill \textbf{Schlüsselwörter}: Power BI, Tableau, Looker Studio, ETL, KPI-Reporting}
\cvline{Datenbanken und Data Engineering}{Erfahren \hfill \textbf{Schlüsselwörter}: PostgreSQL, dbt, Datenmodellierung, SQL, Data Warehousing}
\cvline{Kommunikation und Präsentation}{Erfahren \hfill \textbf{Schlüsselwörter}: Stakeholder-Management, Storytelling, Workshops, Dokumentation}
\section{Auszeichnungen}

\cventry{Juni 2023}
        {DataViz Berlin Meetup}
        {Best Data Visualization Award}
        {}
        {}
        {Auszeichnung für ein interaktives Dashboard zur Analyse von Kundenbindung im Einzelhandel, das auf dem jährlichen Community-Event präsentiert wurde.}
\section{Zertifikate}

\cventry{Jan. 2021}
        {Google}
        {Google Data Analytics Professional Certificate}
        {\url{https://www.coursera.org/professional-certificates/google-data-analytics}}
        {}
        {}

\cventry{Aug. 2022}
        {Microsoft}
        {Microsoft Certified: Power BI Data Analyst Associate}
        {\url{https://learn.microsoft.com/en-us/credentials/certifications/data-analyst-associate/}}
        {}
        {}
\section{Projekte}

\cventry{Jan. 2023–März 2023}
        {Interaktives Dashboard zur automatisierten Kundensegmentierung auf Basis von RFM-Analysen.}
        {Kundensegmentierungs-Dashboard}
        {\url{https://github.com/lena-hoffmann/kundensegmentierung}}
        {}
        {\begin{itemize}
\item Entwicklung eines Power-BI-Dashboards zur Visualisierung von Kundensegmenten und Kaufverhalten
\item Automatisierte Datenaufbereitung mit Python und dbt für ein skalierbares Datenmodell
\item Identifikation von vier profitablen Kundensegmenten, die direkt in Marketingkampagnen überführt wurden
\end{itemize}
\textbf{Schlüsselwörter}: Power BI, Python, RFM-Analyse, Segmentierung}

\cventry{Apr. 2020–Juli 2020}
        {Prognosemodell für wöchentliche Absatzzahlen eines Online-Shops.}
        {Absatzprognose für E-Commerce}
        {\url{https://github.com/lena-hoffmann/absatzprognose}}
        {}
        {\begin{itemize}
\item Bachelorarbeit zur Erstellung von Zeitreihenprognosen mit ARIMA- und linearen Regressionsmodellen
\item Datenbereinigung und Feature-Engineering mit Python und Pandas
\item Vergleich der Modellgüte anhand von MAPE und Visualisierung der Ergebnisse in Tableau
\end{itemize}
\textbf{Schlüsselwörter}: Python, Pandas, Zeitreihenanalyse, ARIMA, Tableau}
\section{Interessen}

\cvline{Datenvisualisierung}{Dashboard-Design, Storytelling mit Daten, Tableau Public}
\cvline{Sport}{Laufen, Bouldern, Radfahren}
\section{Ehrenamtliche Tätigkeiten}

\cventry{Jan. 2023–Heute}
        {Ehrenamtliche Datenanalystin}
        {Data for Good Berlin}
        {\url{https://www.dataforgood.berlin}}
        {}
        {\begin{itemize}
\item Unterstützung gemeinnütziger Organisationen bei der Auswertung von Umfragedaten
\item Erstellung eines Excel-basierten Dashboards zur Ressourcenplanung für eine Beratungsstelle
\item Durchführung von Workshops zu Datenqualität und einfacher Visualisierung
\end{itemize}}

\end{document}